\clearpage\thispagestyle{empty}
\pdfbookmark[0]{Quellen zu Hardys Pendeluhr (26. Juni 1826)}{hardy1826-sources}
\begin{center}\textbf{Quellen zu Hardys Pendeluhr (26. Juni 1826)}\end{center}
\noindent Diese KI-Übertragung ins heutige Deutsch folgt dem Bericht \textit{Hardy’s Pendeluhr} in \href{https://gdz.sub.uni-goettingen.de/id/PPN236007467}{den \textit{Werken}, Band~VI} (Göttingen, 1874), S.~453--454. Der Bericht erschien ohne Unterschrift unter \textit{Göttingen} in den \textit{Göttingischen gelehrten Anzeigen}, Nr.~101 (26. Juni 1826), \href{https://api.digitale-sammlungen.de/iiif/image/v2/bsb10538704_00329/full/1550,/0/default.jpg}{S.~1001} und \href{https://api.digitale-sammlungen.de/iiif/image/v2/bsb10538704_00330/full/1550,/0/default.jpg}{S.~1002}. Die Links führen zum Exemplar der Bayerischen Staatsbibliothek.
\par\smallskip\noindent Der Erstdruck bringt die 21 Beobachtungen in einer Tabelle und liest \textit{retardirt}, wo in der Werkausgabe \textit{retradirt} steht. Die Übertragung gibt den beschriebenen Rückgang des täglichen Gangs wieder und folgt den Tabellen der Werkausgabe, einschließlich der beiden hochgestellten Sekundenzeichen und der Leerstelle für den täglichen Gang am 17. Juni. KI-Übertragung und KI-Satz.
\clearpage
